\section*{第 \ref{ch:g5:stars-night-sky} 章——恒星与夜空}

\begin{solution}{exo:g5:stars-night-sky:1}
\omterm{def:g4:solar-system:star}{恒星}就是一个太阳：一个巨大的、炽热发光的球体，自己发出光亮。
它们看起来像小火花，只是因为距离遥远——
即便是离我们最近的一颗，它的光也要走上好几年。
\end{solution}

\begin{solution}{exo:g5:stars-night-sky:2}
是的——它们在正午和午夜一样闪耀。白天的天空
被\omterm{def:g2:air-around-us:air}{空气}散射的阳光照得一片明亮，
小火花在这样的强光旁边根本看不见，就像篝火旁边看不见一支蜡烛。
\end{solution}

\begin{solution}{exo:g5:stars-night-sky:3}
蓝白色的那颗——\omterm{def:g4:solar-system:star}{恒星}颜色由橙色（最冷）
经黄色，到蓝白色（最热）。
发橙光的煤与烧至白热的钢，讲的是同一个道理。
\end{solution}

\begin{solution}{exo:g5:stars-night-sky:4}
\omterm{def:g5:stars-night-sky:constellation}{星座}是由几颗亮星组成的有名图形，是为了辨认天空、帮助记忆而人为设想的。
它的星星并不是真正的近邻，
只是从地球上看去恰好排成一线。
\end{solution}

\begin{solution}{exo:g5:stars-night-sky:5}
先找到北斗七星；取斗身前端的两颗指极星，
沿它们的连线在天空中向外延伸约五个间距；
所到达的那颗星便是\omterm{met:g5:stars-night-sky:northstar}{北极星}，
\omterm{met:g5:stars-night-sky:northstar}{北极星}下方的地平线便是北方。
\end{solution}

\begin{solution}{exo:g5:stars-night-sky:6}
它几乎正好位于地球自转轴的上方，
因此天空的旋转几乎不能使它移动——
它是一个固定的北向标志，可代替\omterm{def:g4:magnets-and-compass:compass}{指南针}。
\end{solution}

\begin{solution}{exo:g5:stars-night-sky:7}
地球在自转，每天带着我们转一圈。
同样的自转也使太阳在白天的天空中“移动”——
一次自转，带来两种“移动”。
\end{solution}

\begin{solution}{exo:g5:stars-night-sky:8}
无数颗暗淡的星星，密集得肉眼无法分辨：
这就是\omterm{ex:g5:stars-night-sky:milkyway}{银河}——我们身处的扁平星群，从内部侧看过去的样子。
\end{solution}

\begin{solution}{exo:g5:stars-night-sky:9}
旋转的地球使整个天空绕着穿过\omterm{met:g5:stars-night-sky:northstar}{北极星}的轴转动。
靠近这一静止点的星星沿小圆周运动——相机记录下的是圆环。
朝东的方向上，旋转使星星沿倾斜路径升出地平线——
相机记录下的是上升的光迹。
同一次自转，两种景象。
\end{solution}

\begin{solution}{exo:g5:stars-night-sky:10}
方向：北斗的指极星为她指出\omterm{met:g5:stars-night-sky:northstar}{北极星}，
进而指出北方——无需\omterm{def:g4:magnets-and-compass:compass}{指南针}。
时间：北斗七星本身绕\omterm{met:g5:stars-night-sky:northstar}{北极星}转动，好似一只大时钟的指针，
每天转一整圈；
看它转过多少角度，就知道夜晚过去了多久。
\end{solution}

\begin{solution}{exo:g5:stars-night-sky:11}
不是。你看到的是\omterm{def:g4:solar-system:star}{恒星}\emph{过去}的样子，
即几年前、今晚这束光离开它时的样子——信使，而非发送者。
因此，每台望远镜看到的都是过去：\omterm{def:g4:solar-system:star}{恒星}越远，消息越旧——
这正是遥远的光
对天文学家如此珍贵的原因。
\end{solution}
